% Options for packages loaded elsewhere
\PassOptionsToPackage{unicode}{hyperref}
\PassOptionsToPackage{hyphens}{url}
\PassOptionsToPackage{dvipsnames,svgnames,x11names}{xcolor}
%
\documentclass[
  letterpaper,
  DIV=11,
  numbers=noendperiod]{scrartcl}

\usepackage{amsmath,amssymb}
\usepackage{lmodern}
\usepackage{iftex}
\ifPDFTeX
  \usepackage[T1]{fontenc}
  \usepackage[utf8]{inputenc}
  \usepackage{textcomp} % provide euro and other symbols
\else % if luatex or xetex
  \usepackage{unicode-math}
  \defaultfontfeatures{Scale=MatchLowercase}
  \defaultfontfeatures[\rmfamily]{Ligatures=TeX,Scale=1}
  \setmainfont[]{Times New Roman}
  \setmonofont[]{Times New Roman}
\fi
% Use upquote if available, for straight quotes in verbatim environments
\IfFileExists{upquote.sty}{\usepackage{upquote}}{}
\IfFileExists{microtype.sty}{% use microtype if available
  \usepackage[]{microtype}
  \UseMicrotypeSet[protrusion]{basicmath} % disable protrusion for tt fonts
}{}
\makeatletter
\@ifundefined{KOMAClassName}{% if non-KOMA class
  \IfFileExists{parskip.sty}{%
    \usepackage{parskip}
  }{% else
    \setlength{\parindent}{0pt}
    \setlength{\parskip}{6pt plus 2pt minus 1pt}}
}{% if KOMA class
  \KOMAoptions{parskip=half}}
\makeatother
\usepackage{xcolor}
\usepackage[top=30mm,left=30mm]{geometry}
\setlength{\emergencystretch}{3em} % prevent overfull lines
\setcounter{secnumdepth}{2}
% Make \paragraph and \subparagraph free-standing
\ifx\paragraph\undefined\else
  \let\oldparagraph\paragraph
  \renewcommand{\paragraph}[1]{\oldparagraph{#1}\mbox{}}
\fi
\ifx\subparagraph\undefined\else
  \let\oldsubparagraph\subparagraph
  \renewcommand{\subparagraph}[1]{\oldsubparagraph{#1}\mbox{}}
\fi


\providecommand{\tightlist}{%
  \setlength{\itemsep}{0pt}\setlength{\parskip}{0pt}}\usepackage{longtable,booktabs,array}
\usepackage{calc} % for calculating minipage widths
% Correct order of tables after \paragraph or \subparagraph
\usepackage{etoolbox}
\makeatletter
\patchcmd\longtable{\par}{\if@noskipsec\mbox{}\fi\par}{}{}
\makeatother
% Allow footnotes in longtable head/foot
\IfFileExists{footnotehyper.sty}{\usepackage{footnotehyper}}{\usepackage{footnote}}
\makesavenoteenv{longtable}
\usepackage{graphicx}
\makeatletter
\def\maxwidth{\ifdim\Gin@nat@width>\linewidth\linewidth\else\Gin@nat@width\fi}
\def\maxheight{\ifdim\Gin@nat@height>\textheight\textheight\else\Gin@nat@height\fi}
\makeatother
% Scale images if necessary, so that they will not overflow the page
% margins by default, and it is still possible to overwrite the defaults
% using explicit options in \includegraphics[width, height, ...]{}
\setkeys{Gin}{width=\maxwidth,height=\maxheight,keepaspectratio}
% Set default figure placement to htbp
\makeatletter
\def\fps@figure{htbp}
\makeatother

\usepackage{booktabs}
\usepackage{longtable}
\usepackage{array}
\usepackage{multirow}
\usepackage{wrapfig}
\usepackage{float}
\usepackage{colortbl}
\usepackage{pdflscape}
\usepackage{tabu}
\usepackage{threeparttable}
\usepackage{threeparttablex}
\usepackage[normalem]{ulem}
\usepackage{makecell}
\usepackage{xcolor}
\usepackage{amsmath}
\usepackage{caption}
\KOMAoption{captions}{tableheading}
\makeatletter
\makeatother
\makeatletter
\makeatother
\makeatletter
\@ifpackageloaded{caption}{}{\usepackage{caption}}
\AtBeginDocument{%
\ifdefined\contentsname
  \renewcommand*\contentsname{Table of contents}
\else
  \newcommand\contentsname{Table of contents}
\fi
\ifdefined\listfigurename
  \renewcommand*\listfigurename{List of Figures}
\else
  \newcommand\listfigurename{List of Figures}
\fi
\ifdefined\listtablename
  \renewcommand*\listtablename{List of Tables}
\else
  \newcommand\listtablename{List of Tables}
\fi
\ifdefined\figurename
  \renewcommand*\figurename{Figure}
\else
  \newcommand\figurename{Figure}
\fi
\ifdefined\tablename
  \renewcommand*\tablename{Table}
\else
  \newcommand\tablename{Table}
\fi
}
\@ifpackageloaded{float}{}{\usepackage{float}}
\floatstyle{ruled}
\@ifundefined{c@chapter}{\newfloat{codelisting}{h}{lop}}{\newfloat{codelisting}{h}{lop}[chapter]}
\floatname{codelisting}{Listing}
\newcommand*\listoflistings{\listof{codelisting}{List of Listings}}
\makeatother
\makeatletter
\@ifpackageloaded{caption}{}{\usepackage{caption}}
\@ifpackageloaded{subcaption}{}{\usepackage{subcaption}}
\makeatother
\makeatletter
\makeatother
\ifLuaTeX
  \usepackage{selnolig}  % disable illegal ligatures
\fi
\IfFileExists{bookmark.sty}{\usepackage{bookmark}}{\usepackage{hyperref}}
\IfFileExists{xurl.sty}{\usepackage{xurl}}{} % add URL line breaks if available
\urlstyle{same} % disable monospaced font for URLs
\hypersetup{
  pdftitle={Soft Power in World Politics},
  pdfauthor={Kadir Jun Ayhan},
  colorlinks=true,
  linkcolor={blue},
  filecolor={Maroon},
  citecolor={Blue},
  urlcolor={Blue},
  pdfcreator={LaTeX via pandoc}}

\title{Soft Power in World Politics}
\author{Kadir Jun Ayhan}
\date{1/15/23}

\begin{document}
\maketitle
\renewcommand*\contentsname{Contents}
{
\hypersetup{linkcolor=}
\setcounter{tocdepth}{3}
\tableofcontents
}
\hypertarget{syllabus-2023-spring}{%
\section*{Syllabus (2023 Spring)}\label{syllabus-2023-spring}}

\begin{longtable}[]{@{}
  >{\raggedright\arraybackslash}p{(\columnwidth - 4\tabcolsep) * \real{0.3056}}
  >{\raggedright\arraybackslash}p{(\columnwidth - 4\tabcolsep) * \real{0.4306}}
  >{\raggedright\arraybackslash}p{(\columnwidth - 4\tabcolsep) * \real{0.2639}}@{}}
\toprule()
\endhead
{Course Title} & Soft Power in World Politics & {Course No.:} IS639 \\
{Department/ Major} & GSIS & {Credit/ Hours:} 3 \\
{Class time/ Classroom} & Thursday 15:30 -- 18:15 / IEB 902 & \\
{Instructor} & {Name:} Kadir Jun Ayhan & {Department:} GSIS \\
& {E-mail:} \href{mailto:ayhan@ewha.ac.kr}{\nolinkurl{ayhan@ewha.ac.kr}}
& {Phone:} 02-3277-4628 \\
{Office Hours/ Office Location} &
\href{https://calendly.com/kadirayhan/officehours}{Make an appointment
here} & \\
\bottomrule()
\end{longtable}

\hypertarget{course-overview}{%
\section{\texorpdfstring{{\textbf{Course
Overview}}}{Course Overview}}\label{course-overview}}

\hypertarget{course-description}{%
\subsection{\texorpdfstring{\textbf{Course
Description}}{Course Description}}\label{course-description}}

The main target of this class is graduate students who are interested in
power dynamics in world politics. The course starts off with definitions
of power and its different forms. We will learn about how power
functions in world politics with examples from global governance
issue-areas, including direct and diffuse kinds of power. We will
critically approach the concepts of soft power and smart power with the
questions ``how soft is soft power'' and ``can there be pure persuasion
in world politics'' in mind. Furthermore, we will discuss measurement of
(soft) power and influence in world politics. The course follows an
active learning pedagogy, aiming to make students take learning into
their own hands.

\hypertarget{prerequisites}{%
\subsection{\texorpdfstring{\textbf{Prerequisites}}{Prerequisites}}\label{prerequisites}}

Students are expected to have at least introduction-level knowledge of
international relations theories.

\hypertarget{course-format}{%
\subsection{\texorpdfstring{\textbf{Course
Format}}{Course Format}}\label{course-format}}

\begin{longtable}[]{@{}
  >{\raggedright\arraybackslash}p{(\columnwidth - 4\tabcolsep) * \real{0.2603}}
  >{\raggedright\arraybackslash}p{(\columnwidth - 4\tabcolsep) * \real{0.3836}}
  >{\raggedright\arraybackslash}p{(\columnwidth - 4\tabcolsep) * \real{0.3562}}@{}}
\toprule()
\endhead
{Lecture} & {Discussion/ Presentation} & {Experiment/ Practicum} \\
20\% & 40\% & 40\% \\
\bottomrule()
\end{longtable}

This course is designed to facilitate student-centric learning. The
assignments for this course are not merely for grading but designed in a
way to help students practice and improve their critical thinking,
researching, presentation and academic writing skills. There are no
in-class exams for this course. This is because, I believe that you put
so much effort in writing papers in exams to be read only by one person.
On the other hand, writing papers give students an opportunity to
research and keep their writings as working papers which can be improved
in the future.

Active participation in discussion is strongly encouraged for this
course. In order to motivate students to read all required articles and
come to class prepared for discussion, response papers and discussion
participation constitute 25\% of students' grade for this course. In the
beginning of every lecture, you will be asked to share what you have
written in your response papers to stimulate discussion (based on what
you have already prepared). We will discuss your response papers in the
first part of the class. The second part of the class will begin with my
summary and clarification of some of the arguments and concepts in the
assigned readings. Then, we will have one or two student presentations
(depending on the number of students) which will apply the theoretical
concepts in the readings to a case study. The requirements of the
assignments are detailed below.

I expect students to attend classes, to arrive on time, and to come to
class prepared to engage in class discussion by reading at least the
required readings, taking note of the key arguments, and identifying
strengths, weaknesses, and contradictions in these arguments.

\hypertarget{course-objectives}{%
\subsection{\texorpdfstring{\textbf{Course
Objectives}}{Course Objectives}}\label{course-objectives}}

Learning objectives:

- Recognizing various theoretical approaches to political power in world
politics.

- Being able to classify, exemplify and interpret different kinds of
political power in world politics.

- Applying theoretical understanding of political power to case studies
(presentation and workshop), and to other new settings (film); and
analyzing and critiquing these cases.

\hypertarget{evaluation-system}{%
\subsection{\texorpdfstring{\textbf{Evaluation
System}}{Evaluation System}}\label{evaluation-system}}

\begin{longtable}[]{@{}
  >{\raggedright\arraybackslash}p{(\columnwidth - 8\tabcolsep) * \real{0.2000}}
  >{\raggedright\arraybackslash}p{(\columnwidth - 8\tabcolsep) * \real{0.2000}}
  >{\raggedright\arraybackslash}p{(\columnwidth - 8\tabcolsep) * \real{0.2000}}
  >{\raggedright\arraybackslash}p{(\columnwidth - 8\tabcolsep) * \real{0.2000}}
  >{\raggedright\arraybackslash}p{(\columnwidth - 8\tabcolsep) * \real{0.2000}}@{}}
\toprule()
\begin{minipage}[b]{\linewidth}\raggedright
{Mid-term Paper}
\end{minipage} & \begin{minipage}[b]{\linewidth}\raggedright
{Final Paper}
\end{minipage} & \begin{minipage}[b]{\linewidth}\raggedright
{Response Papers}
\end{minipage} & \begin{minipage}[b]{\linewidth}\raggedright
{Presentation}
\end{minipage} & \begin{minipage}[b]{\linewidth}\raggedright
{Participation and Discussion}
\end{minipage} \\
\midrule()
\endhead
30\% & 30\% & 10\% & 15\% & 15\% \\
\bottomrule()
\end{longtable}

Explanation of the evaluation system:

\hypertarget{response-papers}{%
\subsubsection{\texorpdfstring{\textbf{\emph{Response
Papers}}}{Response Papers}}\label{response-papers}}

Length: Less than half a page.

In order to prepare for class discussion, you must read all required
articles and submit weekly response papers beginning from Week 4. Each
student needs to submit at least four response papers (that is for four
different weeks). The presenter(s) should not submit a response paper.
For the response papers, students will find a recent (last six months)
news article (from newspapers, magazines or semi-academic journals) and
connect it to the required readings for that week (critically analyzing
whether it is in line with the arguments in the required readings or
not). Make sure to copy-paste the news article as well as its link below
your response paper. Weekly response papers need to be submitted to
Cyber Campus before class time. You are encouraged to, but you don't
have to, make comments to your colleagues' response papers.

\hypertarget{presentations}{%
\subsubsection{\texorpdfstring{\textbf{\emph{Presentations}}}{Presentations}}\label{presentations}}

Length: Around 15 minutes

All students will make one presentation on a case study. The presenter
must read the assigned readings carefully when preparing for the
presentation. You are asked to present a case study that is related to
the required readings. You must make use of the theories/ frameworks/
concepts in the assigned readings for your analysis of the case study of
your choice. For example, if the required article looks at power
dynamics of the US-led international liberal order, you can apply a
similar approach to look at Soviet Union-led regional international
order. Make sure that you are not merely summarizing another article.
The point of the presentation is for you to apply what you learned from
the readings to a case study. You must show your analysis and your
value-added. Discussion will follow presentations (based on discussion
questions prepared by the presenter), and all students must actively
participate in this discussion.

\hypertarget{sec-workshop}{%
\subsubsection{\texorpdfstring{\textbf{\emph{Workshop and Mid-term
Paper}}}{Workshop and Mid-term Paper}}\label{sec-workshop}}

Due date: 20 Apr 23:59

The mid-term paper will be based on the workshop when students will
discuss how we can measure power/ influence in world politics in groups.
Within the next week, the same groups will finalize their team report
(1000 words) based on their discussions and further research and submit
it to Cyber Campus (Turnitin).

The reports will be presented on 27 Apr followed by debriefing.

Please note that there will not be a class on 20 Apr in lieu of mid-term
exam week.

\hypertarget{sec-final}{%
\subsubsection{\texorpdfstring{\textbf{\emph{Final
Paper}}}{Final Paper}}\label{sec-final}}

Due date: 15 Jun 23:59

Students will watch a film (to be confirmed later), and write an essay
(1600-2000 words -- excluding references), identifying different kinds
of power relations in the film based on what we covered in this class
(in particular in reference to Barnett and Duvall's power typology). I
will upload the guidelines for the essay on Cyber Campus.

\hypertarget{course-materials-and-additional-readings}{%
\section{\texorpdfstring{{\textbf{Course Materials and Additional
Readings}}}{Course Materials and Additional Readings}}\label{course-materials-and-additional-readings}}

\hypertarget{required-materials}{%
\subsection{Required Materials}\label{required-materials}}

You can borrow/ rent/ or buy the following books for this course since
we use one or two chapters from them, but otherwise there are no
required textbooks:

Barnett, M., \& Duvall, R. (2005). Power in Global Governance. Cambridge
University Press.

Berenskoetter, F., \& Williams, M. J. (2007). Power in World Politics.
Routledge.

Nye, J. S. (2011). The Future of Power. Public Affairs.

\hypertarget{supplementary-materials}{%
\subsection{Supplementary Materials}\label{supplementary-materials}}

Refer to the detailed course schedule below.

\hypertarget{optional-additional-readings}{%
\subsection{Optional Additional
Readings}\label{optional-additional-readings}}

Refer to the detailed course schedule below.

\hypertarget{course-policies}{%
\section{\texorpdfstring{{\textbf{Ⅲ. Course
Policies}}}{Ⅲ. Course Policies}}\label{course-policies}}

\hypertarget{citation-and-plagiarism}{%
\subsection{Citation and Plagiarism}\label{citation-and-plagiarism}}

Students must use proper citation and avoid plagiarism. Any citation
style is fine as long as it is consistent throughout the paper.
Plagiarism will not be tolerated and severely punished. The papers will
be uploaded to Turnitin (via Cyber Campus).

\href{(https://www.nyu.edu/about/policies-guidelines-compliance/policies-and-guidelines/academic-integrity-for-students-at-nyu.html)}{``Plagiarism:
presenting others' work without adequate acknowledgement of its source,
as though it were one's own. Plagiarism is a form of fraud. We all stand
on the shoulders of others, and we must give credit to the creators of
the works that we incorporate into products that we call our own. Some
examples of plagiarism:}

\begin{itemize}
\item
  \href{(https://www.nyu.edu/about/policies-guidelines-compliance/policies-and-guidelines/academic-integrity-for-students-at-nyu.html)}{a
  sequence of words incorporated without quotation marks;}
\item
  \href{(https://www.nyu.edu/about/policies-guidelines-compliance/policies-and-guidelines/academic-integrity-for-students-at-nyu.html)}{an
  unacknowledged passage paraphrased from another's word;}
\item
  \href{(https://www.nyu.edu/about/policies-guidelines-compliance/policies-and-guidelines/academic-integrity-for-students-at-nyu.html)}{the
  use of ideas, sound recordings, computer data or images created by
  others as though it were one's own.''}
\end{itemize}

See
\href{https://owl.purdue.edu/owl/research_and_citation/resources.html}{also
this link}.

\hypertarget{late-submissions}{%
\subsection{Late Submissions}\label{late-submissions}}

For late submissions, you will get 80\% of your grading unless you have
a valid excuse. You can always contact me if you have a valid excuse to
ask for extension. I do not require students to submit official
documents (doctor report etc.). Your word is enough for me.

\newpage{}

\hypertarget{course-schedule}{%
\section{\texorpdfstring{{Course
Schedule}}{Course Schedule}}\label{course-schedule}}

\begin{longtable}{rll}
\toprule
Week & Date & Topics \& Class Materials, Assignments \\ 
\midrule
1 & 02 Mar & Introduction of the Course \\ 
2 & 09 Mar & The Concept of Political Power \\ 
3 & 16 Mar & No class - (Make-up class later) \\ 
4 & 23 Mar & The Concepts of Soft Power and Smart Power \\ 
5 & 30 Mar & Skeptical Views on Soft Power \\ 
6 & 06 Apr & Measuring (Soft) Power \\ 
7 & 13 Apr & Workshop: What Do Soft Power Indices Actually Measure? \\ 
8 & 20 Apr & Mid-term Week \\ 
9 & 27 Apr & Mid-term Paper Presentations and Debriefing \\ 
10 & 04 May & Legitimacy and Power \\ 
11 & 11 May & Reputation, Status and Power \\ 
12 & 18 May & Compulsory Power \\ 
13 & 25 May & Institutional Power \\ 
14 & 01 Jun & Structural Power \\ 
15 & 08 Jun & Productive Power \\ 
16 & 15 Jun & Submission of Final Papers \\ 
17 & 22 Jun &  \\ 
\bottomrule
\end{longtable}

\newpage{}

\hypertarget{detailed-course-schedule-with-readings}{%
\section{\texorpdfstring{{Detailed Course Schedule with
Readings}}{Detailed Course Schedule with Readings}}\label{detailed-course-schedule-with-readings}}

\hypertarget{week-1-introduction-of-the-course}{%
\subsection*{Week 1: Introduction of the
Course}\label{week-1-introduction-of-the-course}}
\addcontentsline{toc}{subsection}{Week 1: Introduction of the Course}

Introduction of the course including course contents, assignments, and
expectations.

\hypertarget{week-2-the-concept-of-political-power}{%
\subsection*{Week 2: The Concept of Political
Power}\label{week-2-the-concept-of-political-power}}
\addcontentsline{toc}{subsection}{Week 2: The Concept of Political
Power}

\textbf{Required Readings:}

\begin{itemize}
\tightlist
\item
  Baldwin, D. A. (2021). The faces of power revisited. \emph{Journal of
  Political Power}, \emph{14}(1), 85--96.
  \url{https://doi.org/10.1080/2158379X.2021.1878407}
\item
  Barnett, M., \& Duvall, R. (2005). Power in international politics.
  \emph{International Organization}, \emph{59}(1), 39--75.
  \url{https://doi.org/10.1017/S0020818305050010}
\item
  Berenskoetter, F. (2007). Thinking about power. In F. Berenskoetter \&
  M. J. Williams (Eds.), \emph{Power in world politics} (pp. 1--22).
  Routledge. \url{https://doi.org/10.4324/9780203944691}
\end{itemize}

\textbf{Recommended Readings:} \{.unnumbered\}

\begin{itemize}
\tightlist
\item
  Baldwin, D. A. (2012). Power and International Relations. Handbook of
  International Relations. Walter Carlsnaes, T. Risse and B. A. Simmons.
  London, UK, SAGE.
\item
  Avant, D. D., Finnemore, M., \& Sell, S. K. (2010). Who Governs the
  Globe? In D. D. Avant, M. Finnemore, \& S. K. Sell (Eds.), Who Governs
  the Globe?. Cambridge, UK: Cambridge University Press.
\item
  Murphy, J. B. (2011). ``Perspectives on power.'' Journal of Political
  Power 4(1): 87-103.
\end{itemize}

\hypertarget{week-3-no-class---make-up-class-later}{%
\subsection*{Week 3: No class - (Make-up class
later)}\label{week-3-no-class---make-up-class-later}}
\addcontentsline{toc}{subsection}{Week 3: No class - (Make-up class
later)}

\hypertarget{week-4-the-concepts-of-soft-power-and-smart-power}{%
\subsection*{Week 4: The Concepts of Soft Power and Smart
Power}\label{week-4-the-concepts-of-soft-power-and-smart-power}}
\addcontentsline{toc}{subsection}{Week 4: The Concepts of Soft Power and
Smart Power}

\textbf{Required Readings:}

\begin{itemize}
\tightlist
\item
  Nye, J. S. (2011). \emph{The future of power}. Public Affairs.
  \url{https://www.publicaffairsbooks.com/titles/joseph-s-nye/the-future-of-power/9781586488925/},
  Chapters 4 and 7.
\end{itemize}

\textbf{Recommended Readings:}

\begin{itemize}
\item
  Armitage, R. L., \& Nye, J. S. (2007). CSIS Commission on Smart Power:
  A Smarter, More Secure America: CSIS. Vuving, A. (2009). How soft
  power works. American Political Science Association Annual Meeting.
  Toronto.
\item
  Hayden, C. (2012). The Rhetoric of Soft Power. Lexington Books.
  Ohnesorge, H. W. (2020). Soft Power: The Forces of Attraction in
  International Relations. Springer.
\item
  Lee, G. (2009). A theory of soft power and Korea's soft power
  strategy. Korean Journal of Defense Analysis, 21(2), 205-218.
  https://doi.org/10.1080/10163270902913962
\item
  Nye, J. S. (2004). Soft Power: The Means to Success in World Politics.
  New York, Public Affairs.
\item
  Nye, J. S. (2018). How Sharp Power Threatens Soft Power: The Right and
  Wrong Ways to Respond to Authoritarian Influence. Foreign Affairs.
  https://www.foreignaffairs.com/articles/china/2018-01-24/how-sharp-power-threatens-soft-power
\item
  Nye, J. S. (2019). Soft Power and Public Diplomacy Revisited. The
  Hague Journal of Diplomacy, 14(1-2), 7-20.
  https://doi.org/10.1163/1871191x-14101013
\item
  Nye, J. S. (2021). Soft power: the evolution of a concept. Journal of
  Political Power, 14(1), 196-208.
  https://doi.org/10.1080/2158379X.2021.1879572
\item
  Van Ham, P. (2010). Social Power in International Politics. New York,
  Routledge. Lukes, S. (2005). ``Power and the Battle for Hearts and
  Minds.'' Millennium 33(3): 477-493.
\item
  Walker, C., \& Ludwig, J. (Eds.). (2017). Sharp Power: Rising
  Authoritarian Influence. National Endowment for Democracy.
\end{itemize}

\hypertarget{week-5-skeptical-views-on-soft-power}{%
\subsection*{Week 5: Skeptical Views on Soft
Power}\label{week-5-skeptical-views-on-soft-power}}
\addcontentsline{toc}{subsection}{Week 5: Skeptical Views on Soft Power}

\textbf{Required Readings:}

\begin{itemize}
\tightlist
\item
  Bially Mattern, J. (2005). Why {``soft power''} isn't so soft:
  Representational force and the sociolinguistic construction of
  attraction in world politics. \emph{Millennium - Journal of
  International Studies}, \emph{33}(3), 583--612.
  \url{https://doi.org/10.1177/03058298050330031601}
\item
  Kearn, D. W. (2011). The hard truths about soft power. \emph{Journal
  of Political Power}, \emph{4}(1), 65--85.
  \url{https://doi.org/10.1080/2158379x.2011.556869}
\end{itemize}

\textbf{Recommended Readings:}

\begin{itemize}
\tightlist
\item
  Ayhan, K. J. (2021). Rethinking Soft Power: Instilling Fear,
  Satisfying Appetite or Appealing to Spirit. (Unpublished manuscript)
\item
  Lebow, R. N. (2005). Power, Persuasion and Justice. Millennium, 33(3),
  551-581.
\item
  Manor, I. (2019, October 3). The Banality of Soft Power.
  https://digdipblog.com/2019/10/03/the-banality-of-soft-power/
\end{itemize}

\hypertarget{week-6-measuring-soft-power}{%
\subsection*{Week 6: Measuring (Soft)
Power}\label{week-6-measuring-soft-power}}
\addcontentsline{toc}{subsection}{Week 6: Measuring (Soft) Power}

\textbf{Required Readings:}

\begin{itemize}
\tightlist
\item
  Beckley, M. (2018). The power of nations: Measuring what matters.
  \emph{International Security}, \emph{43}(2), 7--44.
  \url{https://doi.org/10.1162/isec_a_00328}
\item
  Pomeroy, C., \& Beckley, M. (2019). Correspondence: Measuring power in
  international relations. \emph{International Security}, \emph{44}(1),
  197--200. \url{https://doi.org/10.1162/isec_c_00355}
\item
  Yun, S.-H. (2018). An overdue critical look at soft power measurement:
  The construct validity of the soft power 30 in focus. \emph{Journal of
  International and Area Studies}, \emph{25}(2), 1--20.
  \url{https://doi.org/10.2307/26909941}
\end{itemize}

\textbf{Recommended Readings:}

\begin{itemize}
\tightlist
\item
  Barnett, G. A., Xu, W. W., Chu, J., Jiang, K., Huh, C., Park, J. Y.,
  \& Park, H. W. (2017). Measuring international relations in social
  media conversations. Government Information Quarterly, 34(1), 37-44.
  https://doi.org/10.1016/j.giq.2016.12.004
\item
  Hart, J. (1976). Three Approaches to the Measurement of Power in
  International Relations. International Organization, 30(2), 289-305.
  doi:10.1017/S0020818300018282
\item
  Mcclory, J. (Ed.). (2019). Soft Power 30: A Global Ranking of Soft
  Power 2019. Portland.
\item
  Treverton, G., \& Jones, S. G. (2005). Measuring Power. Harvard
  International Review, 27(2), 54-58.
\end{itemize}

\hypertarget{week-7-workshop-what-do-soft-power-indices-actually-measure}{%
\subsection*{Week 7: Workshop: What Do Soft Power Indices Actually
Measure?}\label{week-7-workshop-what-do-soft-power-indices-actually-measure}}
\addcontentsline{toc}{subsection}{Week 7: Workshop: What Do Soft Power
Indices Actually Measure?}

See Workshop and Mid-term Section (Section~\ref{sec-workshop}) for
details.

\hypertarget{week-8-mid-term-week}{%
\subsection*{Week 8: Mid-term Week}\label{week-8-mid-term-week}}
\addcontentsline{toc}{subsection}{Week 8: Mid-term Week}

See Workshop and Mid-term Section (Section~\ref{sec-workshop}) for
details.

\hypertarget{week-9-mid-term-paper-presentations-and-debriefing}{%
\subsection*{Week 9: Mid-term Paper Presentations and
Debriefing}\label{week-9-mid-term-paper-presentations-and-debriefing}}
\addcontentsline{toc}{subsection}{Week 9: Mid-term Paper Presentations
and Debriefing}

See Workshop and Mid-term Section (Section~\ref{sec-workshop}) for
details.

\hypertarget{week-10-legitimacy-and-power}{%
\subsection*{Week 10: Legitimacy and
Power}\label{week-10-legitimacy-and-power}}
\addcontentsline{toc}{subsection}{Week 10: Legitimacy and Power}

\textbf{Required Readings:}

\begin{itemize}
\tightlist
\item
  Finnemore, M. (2009). Legitimacy, hypocrisy, and the social structure
  of unipolarity. \emph{World Politics}, \emph{61}(1), 58--85.
  \url{https://doi.org/10.1017/CBO9780511996337.003}
\item
  Hurd, I. (1999). Legitimacy and authority in international politics.
  \emph{International Organization}, \emph{53}(2), 379--408.
  \url{https://doi.org/10.1162/002081899550913}
\end{itemize}

\textbf{Recommended Readings:}

\begin{itemize}
\tightlist
\item
  Bukovansky, M. (2002). Legitimacy and Power Politics: The American and
  French Revolutions in International Political Culture. Princeton
  University Press.
\item
  Clark, I. (2005). Legitimacy in International Society. Oxford
  University Press.
\item
  Clark, I. (2007). International Legitimacy and World Society. Oxford
  University Press.
\item
  Franck, T. M. (1990). The Power of Legitimacy Among Nations. Oxford
  University Press.
\item
  Franck, T. M. (1995). Fairness in International Law and Institutions.
  Oxford University Press.
\item
  Hall, R. B. (1997). Moral Authority as a Power Resource. International
  Organization, 51(4), 591-622.
\item
  Hurd, I. (2008). After Anarchy: Legitimacy and Power in the United
  Nations Security Council. Princeton University Press.
\item
  Kavalski, E. (2013). The struggle for recognition of normative powers:
  Normative power Europe and normative power China in context.
  Cooperation and Conflict, 48(2), 247-267.
\item
  Mulligan, S. P. (2006). The uses of legitimacy in international
  relations. Millennium, 34(2), 349-375.
\item
  Suchman, M. C. (1995). Managing legitimacy: Strategic and
  institutional approaches. Academy of management Review, 20(3),
  571-610.
\item
  Weber, M. (2019). Economy and Society: A New Translation (K. Tribe,
  Ed. \& Trans.). Harvard University Press.
\end{itemize}

\hypertarget{week-11-reputation-status-and-power}{%
\subsection*{Week 11: Reputation, Status and
Power}\label{week-11-reputation-status-and-power}}
\addcontentsline{toc}{subsection}{Week 11: Reputation, Status and Power}

\textbf{Required Readings:}

\begin{itemize}
\tightlist
\item
  Khong, Y. F. (2019). Power as prestige in world politics.
  \emph{International Affairs}, \emph{95}(1), 119--142.
  \url{https://doi.org/10.1093/ia/iiy245}
\item
  Lebow, R. N. (2005). Power, persuasion and justice. \emph{Millennium},
  \emph{33}(3), 551--581.
  \url{https://doi.org/10.1177/03058298050330031501}
\item
  Wohlforth, W. C., Carvalho, B. de, Leira, H., \& Neumann, I. B.
  (2017). Moral authority and status in international relations: Good
  states and the social dimension of status seeking. \emph{Review of
  International Studies}, 1--21.
  \url{https://doi.org/10.1017/S0260210517000560}
\end{itemize}

\textbf{Recommended Readings:}

\begin{itemize}
\tightlist
\item
  Ayhan, K. J. (2019). ``Rethinking Korea's Middle Power Diplomacy as a
  Nation Branding Project.'' KOREA OBSERVER 50: 1.
\item
  Bry, S. (2015). The Production of Soft Power: Practising Solidarity in
  Brazilian South--South Development Projects. Canadian Journal of
  Development Studies/Revue canadienne d'études du développement, 36(4),
  442-458.
\item
  Brewster, R. (2009). The limits of reputation on compliance.
  International Theory, 1(2), 323-333.
  https://doi.org/10.1017/s175297190999008x
\item
  Cull, N. (2018). The Quest for Reputational Security: The Soft Power
  Agenda of Kazakhstan, USC Center on Public Diplomacy.
\item
  Dafoe, A., et al.~(2014). ``Reputation and Status as Motives for
  War.'' Annual Review of Political Science 17(1): 371-393.
\item
  de Carvalho, B., \& Neumann, I. B. (2014). Small State Status Seeking:
  Norway's Quest for International Standing. Routledge.
\item
  Hwang, K.-D. (2014). Korea's Soft Power as an Alternative Approach to
  Africa in Development Cooperation. African and Asian Studies, 13(3),
  249-271.
\item
  Lebow, R. N. (2008). A Cultural Theory of International Relations.
  Cambridge, Cambridge University Press. Ch. 2
\item
  Renshon, J. (2016). ``Status Deficits and War.'' International
  Organization 70(3): 513-550. Wood, S. (2013). ``Prestige in world
  politics: History, theory, expression.'' International Politics 50(3):
  387-411.
\end{itemize}

\hypertarget{week-12-compulsory-power}{%
\subsection*{Week 12: Compulsory Power}\label{week-12-compulsory-power}}
\addcontentsline{toc}{subsection}{Week 12: Compulsory Power}

Note: Compulsory power through material threats is rather
straightforward and you read articles on these topics in other courses.
Therefore, I will spend more time on compulsory power through naming and
shaming and other symbolic strategies.

\textbf{Required Readings:}

\begin{itemize}
\tightlist
\item
  Murdie, A. M., \& Davis, D. R. (2012). Shaming and blaming: Using
  events data to assess the impact of human rights {INGOs}1.
  \emph{International Studies Quarterly}, \emph{56}(1), 1--16.
  \url{https://doi.org/10.1111/j.1468-2478.2011.00694.x}
\item
  Souva, M. (2022). Material military power: A country-year measure of
  military power, 1865--2019. \emph{Journal of Peace Research},
  00223433221112970. \url{https://doi.org/10.1177/00223433221112970}
\end{itemize}

\textbf{Recommended Readings:}

\begin{itemize}
\tightlist
\item
  Bob, C. (2002). ``Merchants of Morality.'' Foreign Policy(129): 36.
\item
  Keck, M. E. and K. Sikkink (1998). Activists beyond Borders: Advocacy
  Networks in International Politics. Ithaca: NY, Cambridge University
  Press.
\item
  Risse, T., et al., Eds. (1999). The Power of Human Rights
  International Norms and Domestic Change. Cambridge, Cambridge
  University Press.
\item
  Risse, T., et al., Eds. (2013). The Persistent Power of Human Rights:
  From Commitment to Compliance. Cambridge, Cambridge University Press.
\item
  Vadlamannati, K. C., et al.~(2018). ``Human Rights Shaming and FDI:
  Effects of the UN Human Rights Commission and Council.'' World
  Development 104: 222-237.
\end{itemize}

\hypertarget{week-13-institutional-power}{%
\subsection*{Week 13: Institutional
Power}\label{week-13-institutional-power}}
\addcontentsline{toc}{subsection}{Week 13: Institutional Power}

\textbf{Required Readings:}

\begin{itemize}
\tightlist
\item
  Barnett, M., \& Finnemore, M. (2005). The power of liberal
  international organizations. In M. Barnett \& R. Duvall (Eds.),
  \emph{Power in global governance}. Cambridge University Press.
  \url{https://doi.org/10.1017/CBO9780511491207.007}
\item
  Dreher, A., \& Jensen, N. M. (2007). Independent actor or agent? An
  empirical analysis of the impact of u.s. Interests on international
  monetary fund conditions. \emph{The Journal of Law and Economics},
  \emph{50}(1), 105--124. \url{https://doi.org/10.1086/508311}
\item
  Goh, E. (2011). Institutions and the great power bargain in east asia:
  {ASEAN}'s limited {``brokerage''} role. \emph{International Relations
  of the Asia-Pacific}, \emph{11}(3), 373--401.
  \url{https://doi.org/10.1093/irap/lcr014}
\end{itemize}

\textbf{Recommended Readings:}

\begin{itemize}
\tightlist
\item
  Hurrell, A. (2004). Power, institutions, and the production of
  inequality. Power in Global Governance. M. Barnett and R. Duvall.
  Cambridge, Cambridge University Press: 33-58.
\item
  Gronau, J. and H. Schmidtke (2016). ``The quest for legitimacy in
  world politics -- international institutions' legitimation
  strategies.'' Review of International Studies 42(3): 535-557.
\item
  Goldstein, J. and R. O. Keohane, Eds. (1993). Ideas and Foreign
  Policy: An Analytical Framework. Ithaca, NY, Cornell University.
\item
  Ikenberry, G. J. (2001). After Victory: Institutions, Strategic
  Restraint, and the Rebuilding of Order After Major Wars. Princeton
  University Press.
\item
  Ikenberry, G. J. (2019). Reflections on after victory. The British
  Journal of Politics and International Relations, 21(1), 5-19.
\item
  Keohane, R. O. and J. S. Nye (1987). ``Power and Interdependence
  Revisited.'' International Organization 41(4): 725-753.
\item
  Kilby, C. (2009). ``The political economy of conditionality: An
  empirical analysis of World Bank loan disbursements.'' Journal of
  Development Economics 89(1): 51-61.
\item
  Lee, G. (2018). Soft power after Nye: the neoliberal international
  order and soft power learning. Handbook of Cultural Security. Y.
  Watanabe. Cheltenham, Edward Elgar Publishing Limited: 253-268.
\item
  Pouliot, V. (2014). Setting Status in Stone: The Negotiation of
  International Institutional Privileges. In D. Welch Larson, - T.
  Shaffer, G. (2004). Power, governance, and the WTO: a comparative
  institutional approach. Power in Global Governance. M. Barnett and R.
  Duvall. Cambridge, Cambridge University Press: 130-160.
\item
  Paul, V. \& W. C. Wohlforth (Eds.), Status in World Politics
  (pp.~192-216). Cambridge University Press. https://doi.org/DOI:
  10.1017/CBO9781107444409.012
\end{itemize}

\hypertarget{week-14-structural-power}{%
\subsection*{Week 14: Structural Power}\label{week-14-structural-power}}
\addcontentsline{toc}{subsection}{Week 14: Structural Power}

\textbf{Required Readings:}

\begin{itemize}
\tightlist
\item
  Ikenberry, G. J., \& Kupchan, C. A. (1990). Socialization and
  hegemonic power. \emph{International Organization}, \emph{44}(3),
  283--315. \url{https://doi.org/10.1017/S002081830003530X}
\item
  Winecoff, W. K. (2020). {``The persistent myth of lost hegemony,''}
  revisited: Structural power as a complex network phenomenon.
  \emph{European Journal of International Relations}, \emph{26}(1),
  209--252. \url{https://doi.org/10.1177/1354066120952876}
\end{itemize}

\textbf{Recommended Readings:}

\begin{itemize}
\tightlist
\item
  Akins, A. (2012). ````Next Time, Just Remember the Story'' Unlearning
  Empire in Silko's Ceremony.'' Studies in American Indian Literatures
  24(1): 1-14.
\item
  Cox, R. (1981). Social Forces, States and World Order: Beyond
  International Relations Theory. Millennium, 10(2), 126-55.
\item
  Cox, R. (1983). Gramsci, Hegemony and International Relations: An
  Essay in Method. Millennium, 12(2), 162-175.
\item
  Dunch, R. (2002). ``Beyond Cultural Imperialism: Cultural Theory,
  Christian Missions, and Global Modernity.'' History and Theory 41(3):
  301-325.
\item
  Guzzini, S. (1993). Structural power: the limits of neorealist power
  analysis. International Organization, 47(3), 443-478.
\item
  Hobsbawm, E. and T. Ranger, Eds. (1983). The Invention of Tradition.
  Cambridge, Cambridge University Press. Ch. 1
\item
  Lukes, S. (2005). Power: A Radical View. Basingstoke, Palgrave
  Macmillan. Pp. 14-59
\item
  Mann, M. (2004). The first failed empire of the 21st century', Review
  of International Studies, 30(4), 631-653.
\item
  Nexon, D. B. and T. Wright (2007). What's at Stake in the American
  Empire Debate. American Political Science Review, 101(2), 253-271.
\item
  Pamment, J. (2016). Towards a New Conditionality? The Convergence of
  International Development, Nation Brands and Soft Power in the British
  National Security Strategy. Journal of International Relations and
  Development. doi:10.1057/s41268-016-0074-9
\item
  Patomäki, H., \& Teivainen, T. (2004). The World Social Forum: An Open
  Space or a Movement of Movements?. Theory, Culture \& Society, 21 (6),
  45-54
\item
  Peterson, J. H. (2012). ``A Conceptual Unpacking of Hybridity:
  Accounting for Notions of Power, Politics and Progress In Analyses of
  Aid-Driven Interfaces.'' Journal of Peacebuilding \& Development 7(2):
  9-22.
\item
  Petras, J. (1994). ``Cultural Imperialism in Late 20th Century.''
  Economic and Political Weekly 29(32): 2070-2073.
\item
  Pomeranz, K. (2005). ``Empire \& `Civilizing' Missions, Past \&
  Present.'' Daedalus 134(2): 34-45.
\item
  Sautman, B., \& Hairong, Y. (2013). Friends and Interests: China's
  Distinctive Links with Africa. African Studies Review, 50(3), 75-114.
  doi:10.1353/arw.2008.0014
\item
  Strange, S. (1987). The Persistent Myth of Lost Hegemony.
  International Organization, 41(4), 551-574.
\item
  Strange, S. (1990). Finance, information and power. Review of
  International Studies, 16(3), 259-274.
\item
  Worth, O. (2010). Recasting Gramsci in international politics. Review
  of International Studies 37(1), 373-392.
\end{itemize}

\hypertarget{week-15-productive-power}{%
\subsection*{Week 15: Productive Power}\label{week-15-productive-power}}
\addcontentsline{toc}{subsection}{Week 15: Productive Power}

\textbf{Required Readings:}

\begin{itemize}
\tightlist
\item
  Castells, M. (2009). \emph{Communication power}. Oxford University
  Press.
  \url{https://global.oup.com/academic/product/communication-power-9780199681938},
  Chapter 3
\item
  Ipek, P. (2015). Ideas and change in foreign policy instruments: Soft
  power and the case of the turkish international cooperation and
  development agency. \emph{Foreign Policy Analysis}, \emph{11}(2),
  173--193. \url{https://doi.org/10.1111/fpa.12031}
\end{itemize}

\textbf{Recommended Readings:}

\begin{itemize}
\tightlist
\item
  Adler, E., \& Bernstein, S. (2005). Knowledge in Power: The Epistemic
  Construction of Global Governance. In M. Barnett \& R. Duvall (Eds.),
  Power in Global Governance. New York: Cambridge University Press.
\item
  Bae, Y., \& Lee, Y. W. (2020). Socialized soft power: recasting
  analytical path and public diplomacy. Journal of International
  Relations and Development, 23(4), 871-898.
  https://doi.org/10.1057/s41268-019-00169-5
\item
  Bigo, D. (2011). Pierre Bourdieu and International Relations: Power of
  Practices, Practices of Power. International Political Sociology 5(3):
  225--258.
\item
  Bourdieu, P. (1991). Language and Symbolic Power (G. Raymond \& M.
  Adamson, Trans.; J. B. Thompson, Ed.). Polity Press.
\item
  Castells, M. (2015). Communication Power. Oxford: Oxford University
  Press.
\item
  Foucault, M. (1980). Power/ Knowledge (C. Gordon, Trans.). Pantheon.
\item
  Kang, D. C. (2020). International Order in Historical East Asia:
  Tribute and Hierarchy Beyond Sinocentrism and Eurocentrism.
  International Organization, 74(1), 65-93.
\item
  Kinsella, H. M. (2004). Securing the civilian: sex and gender in the
  laws of war. In M. Barnett \& R. Duvall (Eds.), Power in global
  governance (pp.~249-272).
\item
  Merlingen, M. (2003). Governmentality: Towards a Foucauldian Framework
  for the Study of IGOs. Cooperation and Conflict, 38(4): 361-384.
\item
  Morisse-Schilbach, M. (2015). ``Changing the World'': Epistemic
  Communities, and the Democratizing Power of Science. Innovation: The
  European Journal of Social Sciences, 28(1), 18-26.
  doi:10.1080/13511610.2014.943163
\item
  Leander, A. (2005). The Power to Construct International Security: On
  the Significance of Private Military Companies. Millennium, 33(3):
  803-826.
\item
  Onuf, N. (2017). ``The power of metaphor/the metaphor of power.'' The
  Journal of International Communication 23(1): 1-14.
\item
  Schmidt, V. A. (2008). ``Discursive Institutionalism: The Explanatory
  Power of Ideas and Discourse.'' Annual Reviews of Political Science
  11: 303-326.
\item
  Payne, R. A. (2001). ``Persuasion, Frames and Norm Construction.''
  European Journal of International Relations 7(1): 37-61.
\item
  Risse, T. (2011). ``Ideas, discourse, power and the end of the Cold
  War: 20 years on.'' International Politics 48(4-5): 591-606.
\end{itemize}

\hypertarget{week-16-submission-of-final-papers}{%
\subsection*{Week 16: Submission of Final
Papers}\label{week-16-submission-of-final-papers}}
\addcontentsline{toc}{subsection}{Week 16: Submission of Final Papers}

See Final Paper Section (Section~\ref{sec-final}) for details.

\textbf{Recommended Reading:}

\begin{itemize}
\tightlist
\item
  Ayhan, K. J. (2020). Transferring knowledge to narrative worlds:
  Applying power taxonomy to science fiction films. \emph{International
  Studies Perspectives}, \emph{21}(3), 258--274.
  \url{https://doi.org/10.1093/isp/ekz024}
\end{itemize}

\hypertarget{special-accommodations}{%
\section{\texorpdfstring{{Special
Accommodations}}{Special Accommodations}}\label{special-accommodations}}

\begin{itemize}
\tightlist
\item
  According to the University regulation section \#57-3, students with
  disabilities can request for special accommodations related to
  attendance, lectures, assignments, or tests by contacting the course
  professor at the beginning of semester. Based on the nature of the
  students' request, students can receive support for such
  accommodations from the course professor or from the Support Center
  for Students with Disabilities (SCSD). Please refer to the below
  examples of the types of support available in the lectures,
  assignments, and evaluations.
\end{itemize}

\begin{longtable}[]{@{}
  >{\raggedright\arraybackslash}p{(\columnwidth - 4\tabcolsep) * \real{0.3056}}
  >{\raggedright\arraybackslash}p{(\columnwidth - 4\tabcolsep) * \real{0.2083}}
  >{\raggedright\arraybackslash}p{(\columnwidth - 4\tabcolsep) * \real{0.4861}}@{}}
\toprule()
\begin{minipage}[b]{\linewidth}\raggedright
Lecture
\end{minipage} & \begin{minipage}[b]{\linewidth}\raggedright
Assignments
\end{minipage} & \begin{minipage}[b]{\linewidth}\raggedright
Evaluation
\end{minipage} \\
\midrule()
\endhead
Visual impairment: braille, enlarged reading materials. Hearing
impairment: note-taking assistant. Physical impairment: access to
classroom, note-taking assistant. & Extra days for submission,
alternative assignments. & Visual impairment: braille examination paper,
examination with voice support, longer examination hours, note-taking
assistant. Hearing impairment: written examination instead of oral.
Physical impairment: longer examination hours, note-taking assistant. \\
\bottomrule()
\end{longtable}

\begin{itemize}
\item
  Actual support may vary depending on the course.
\item
  If you have other special needs, please let me know. I will do my best
  to flexibly accommodate your needs.
\end{itemize}

\hypertarget{notes}{%
\section{\texorpdfstring{{Notes}}{Notes}}\label{notes}}

The contents of this syllabus are not final. I may update them later.

\hypertarget{read-the-syllabus}{%
\section{\texorpdfstring{{Read the
Syllabus!}}{Read the Syllabus!}}\label{read-the-syllabus}}

\url{https://www.youtube.com/watch?v=xQJ1vQ4RJ5I\&t=1s}

\includegraphics{http://www.phdcomics.com/comics/archive/phd051013s.gif}



\end{document}
